%%%%%%%%%%%%%%%%%%%%%%%%%%%%%%%%%%%%%%%%%
% Compact Academic CV
% LaTeX Template
% Version 2.0 (6/7/2019)
%
% This template originates from:
% https://www.LaTeXTemplates.com
%
% Authors:
% Dario Taraborelli (http://nitens.org/taraborelli/home)
% Vel (vel@LaTeXTemplates.com)
%
% License:
% CC BY-NC-SA 3.0 (http://creativecommons.org/licenses/by-nc-sa/3.0/)
%
%%%%%%%%%%%%%%%%%%%%%%%%%%%%%%%%%%%%%%%%%

%----------------------------------------------------------------------------------------
%	PACKAGES AND OTHER DOCUMENT CONFIGURATIONS
%----------------------------------------------------------------------------------------

\documentclass[11pt]{article} % Default document font size

\input{structure.tex} % Include the file specifying the document structure and styling

% Set PDF meta-information
\hypersetup{
	pdftitle={John Kramer - Curriculum vitae},
	pdfauthor={John Kramer}
}

%----------------------------------------------------------------------------------------

\begin{document}

%----------------------------------------------------------------------------------------
%	CONTACT AND GENERAL INFORMATION
%----------------------------------------------------------------------------------------
\begin{center}
{\LARGE\bfseries John V. Kramer} % Name
\bigskip% Whitespace

Department of Economics, University of Copenhagen, \\ 
Øster Farimagsgade 5, 1353 København, Denmark\\

%\medskip % Whitespace


Email: \href{mailto:john.kramer@econ.ku.dk }{john.kramer@econ.ku.dk }\\ % Email address
Website: \href{https://jvkramer.github.io/}{jvkramer.github.io}\\ % Academic/personal website
\end{center}

%------------------------------------------------

\section*{Post-PhD professional position}

\years{2022-now} Assistant Professor, University of Copenhagen \\ % Current or most recent employment position
\years{2026} Visiting Researcher, Banco de Portugal (Feb-Apr)


\section*{Education}
\years{2022} {Ph.D. in Economics, IIES, Stockholm University} \\
%Thesis committee: Bruno Crépon, Rafael Lalive, Roland Rathelot, Josef Zweimüller \\
%Recommendations:  David Card, Bruno Crépon, Ioana Marinescu, Josef Zweimüller \\
\years{2015} Advanced Studies Program, Kiel Institute for the World Economy \\
\years{2014} M.Sc. in Financial Economics, Maastricht University \\
\years{2012} B.Sc. in Economics and Business Economics, Maastricht University

\section*{Academic visiting}

\years{2018-2019} Visiting student researcher at Stanford University, \\
sponsor: Martin Schneider \\
\years{2010} Visiting student at the University of California---San Diego


% \section*{Areas of specialisation}

% Labor Economics, Applied Econometrics % Primary areas of research interest

%----------------------------------------------------------------------------------------
%	PUBLICATIONS AND TALKS
%----------------------------------------------------------------------------------------

\section*{Research}

\subsection*{Publications}

The Curious Incidence of Monetary Policy Across the Income Distribution, with Tobias Broer and Kurt Mitman, \textbf{\textit{American Economic Journal: Macroeconomics}}, \textit{conditionally accepted} \\

The Distributional Effects of Oil Shocks, with Tobias Broer and Kurt Mitman, \textbf{\textit{IMF Economic Review}}, October 2025 \\

Monetary Policy and Liquidity Constraints: Evidence from the Euro Area, with Mattias Almgren, Jos\'e Gallegos and Ricardo Lima, \textbf{\textit{American Economic Journal: Macroeconomics}}, October 2022 



\subsection*{Selected works in progress}

The Cyclicality of Earnings Growth along the Distribution - Causes and Consequences, WP 2026. \\

It Runs in the Family: Occupational Choice and the Allocation of Talent, with Mattias Almgren and Josef Sigurdsson, R\&R at \textit{The Economic Journal}. \\

Beliefs about the Time Cost of Investing: Experimental Evidence from Denmark, with Christina Gravert, WP 2026. \\

High frequency wage rigidity in Denmark, with Tobias Renkin. \\

The Exchange Rate, with Sebastian Fanelli, Marcus Hagedorn and Kurt Mitman.

\section*{Presentations and discussions (including scheduled)}

\subsection*{Presentations in seminars, workshops, conferences}
\years{2026} \textbf{Invited presentations:} Banco de Portugal

\years{2025} \textbf{Conferences:} New Frontiers for Policy Evaluation (NYU Abu Dhabi), SED conference (Copenhagen), Lisbon Macro Workshop

\years{2024} \textbf{Invited presentations:} Danmarks Nationalbank

\years{2023} \textbf{Conferences:} NORMAC

\years{2022}\textbf{Invited presentations:} Bank for International Settlements; Ludwig Maximilian Universit\"at, Munich; University of Copenhagen; Uppsala University; Queen Mary University; McGill University; University of Lausanne; Nova School of Business and Economics; Copenhagen Business School; University of Oslo; University of Rome Tor Vergata; Federal Reserve Board; Sveriges Riksbank; Danmarks Nationalbank; University of Helsinki; Banque de France.

\textbf{Conferences:} Society of Economics Dynamics; European Economic Association; Oslo Macro Days.

\years{2021} \textbf{Invited presentations:} Stockholm University, NHH Bergen. 

\textbf{Conferences:} European Economic Association. 
 

%----------------------------------------------------------------------------------------
%	TEACHING
%----------------------------------------------------------------------------------------
\section*{Professional activities}
\subsection*{Refereeing services}
\textit{The Review of Economic Studies}, \textit{The Review of Financial Studies}, \textit{The Journal of the European Economic Association}, \textit{The Economic Journal}, \textit{Quantitative Economics}, \textit{The Journal of Public Economics}

\subsection*{Departmental services}
\years{2022-now} Co-organizer of the CBS-University of Copenhagen-Nationalbanken Macro Presentations Copenhagen (MPC)\\
 

%\section*{Affiliation}
%\years{2017-now} Research affiliate at IZA\\

\section*{Research grants, fellowships and awards}

\subsection*{Individual grants}
\years{2026} Banco de Portugal visiting researcher grant \\
\years{2018,2019} Mannerfeldt mobility grant \\
\years{2016} Hedelius Research Fellowship \\
\years{2010,2011} Maastricht University Top 3\% award

\subsection*{Funded research projects}
\years{2025} ``ASKing the right question: Stock market participation in Denmark", funded by the Economic Policy Research Network (EPRN), DKK 495,600 (incl.\ 20\% overhead), with Christina Gravert (PI: Kramer)

%\years{2022-2023} ``The Determinants of Disparities in Reservation Wages", funded by the Upjohn Institute Early Career Research Award with Maxim Massenkoff (PI: Daphne Skandalis)  \\
%\years{2021-2025} ``What is keeping unemployed workers from finding employment? The role of information, job characteristics, and behavioral biases" funded by the Rockwool fundation, joint with Steffen Altmann, Nikolaj Harmon, Robert Mahlstedt, Alexander Sebald (PI: Nikolaj Harmon and Alexander Sebald)  \\
%\years{2021-2022} ``Financial Frictions & the Cost of Job Displacement" funded by the Economic Policy Research Network, joint with Niels Johannesen \\
%\years{2012-2014}  “Counseling of young unemployed job seekers in search clubs”, funded by the European Commission, program progress 2011, joint with Gerard van den Berg, Bruno Crepon, Sylvie Blasco, and Arne Uhlendorff (PI: Gerard van den Berg and Bruno Crepon)


\section*{Teaching}
\years{Fall 2022-} University of Copenhagen : Macroeconomics III (Undergraduate) \\
\years{2018-2019} Stockholm University: Macroeconomics I (graduate)---teaching assistant to Professor Timo Boppart
%\years{2013-2014}	Paris 1 University (TA): Introduction to econometrics 							(Undergraduate), 	Introduction to econometrics using Stata 						(Undergraduate), Public finance (Undergraduate)

%\subsection*{Student supervision}

 %\textbf{Master thesis supervision:} Anna Poulsen (2023); Alberte Hoeg Olsen (2023); Renas Sertdemir (2023); 

%\section*{Policy reports}
%\textbf{Inflation Expectations in Times of COVID-19}, with Olivier Armantier, Gizem Kosar, Rachel Pomerantz, Kyle Smith, Giorgio Topa, and Wilbert van der Klaauw, Liberty Street Economics, May 2020\\
%\textbf{How Widespread Is the Impact of the COVID-19 Outbreak on Consumer Expectations?}, Olivier Armantier, Gizem Kosar, Rachel Pomerantz, Kyle Smith, Giorgio Topa, and Wilbert van der Klaauw, Liberty Street Economics, April 2020\\
%\textbf{Coronavirus Outbreak Sends Consumer Expectations Plummeting}, Olivier Armantier, Gizem Kosar, Rachel Pomerantz, Kyle Smith, Giorgio Topa, and Wilbert van der Klaauw, Liberty Street Economics, April 2020 \\
%\textbf{Did Changes in Economic Expectations Foreshadow Swings in the 2018 Elections?}, with Olivier Armantier, Michael Neubauer, and Wilbert van der Klaauw, Liberty Street Economics, May 2019 \\
%\textbf{Economic Expectations Grow Less Polarized since the 2016 Election}, with Olivier Armantier, Michael Neubauer, and Wilbert van der Klaauw, Liberty Street Economics, Liberty Street Economics, May 2019 \\
%\textbf{L'externalisation des programmes d’accompagnement renforcé des demandeurs d’emploi}, with Bruno Crépon, Etudes et Recherches - Pôle emploi, June 2015\\
%\textbf{Club Jeunes chercheurs d'emploi: Evaluation d'une action pilote}, with Sylvie Blasco, Bruno Crépon, Arne Uhlendorff, Gerard J. van den Berg, Élodie Alberola and Francois Aventur, Etudes et Recherches - Pôle emploi, October 2015\\
%\textbf{Quotient familial, quotient conjugal, impôt individualisé : quels sont les enjeux du débat?} with Sébastien Grobon, Regards croisés sur l'économie 2014/2 (n° 15)





\section*{Miscellaneous}

\years{Language} German (native), English (fluent), Portuguese (beginner), Swedish (beginner), Danish (beginner)	\\
%\years{Software} Stata, Python, SAS, Matlab

%----------------------------------------------------------------------------------------
%	FINAL FOOTER
%----------------------------------------------------------------------------------------

% Any final footer text such as a URL to the latest version of this CV, last updated date, compiled in XeTeX, etc
\begin{center}
	\scriptsize
	Last updated: \today
\end{center}

%----------------------------------------------------------------------------------------

\end{document}
